\begin{exercise}[Transversality of the QED vacuum polarization]{I-CH09-011}
In Euclidean spinor QED, let
the continued loop dimension be $d_{\mathrm{reg}}$ and use
$\{\gamma^\mu,\gamma^\nu\}=2\delta^{\mu\nu}$,
\[
  S(q)=\frac{\slashed q-\ii m_0}{q^2+m_0^2},
  \qquad
  \Pi_{\alpha\beta}(p)=e_0^2\int\frac{\dd^{d_{\mathrm{reg}}}q}{(2\pi)^{d_{\mathrm{reg}}}}
  \operatorname{tr}[\gamma_\alpha S(q)\gamma_\beta S(q+p)].
\]
Contract with $p^\alpha$ and use
\[
  \slashed p=(\slashed q+\slashed p-\ii m_0)
             -(\slashed q-\ii m_0)
\]
to reduce the result to a difference of one-propagator integrals.  Use the
translation invariance of dimensional regularization to prove
$p^\alpha\Pi_{\alpha\beta}=0$.  State why a hard momentum boundary can spoil
this intermediate step and how a gauge-invariant prescription restores the
identity.  Check the result at $p=0$ and deduce the allowed tensor form of
$\Pi_{\alpha\beta}$.
\end{exercise}
